% !TEX program = xelatex
% 공통수학2 · Ⅱ. 집합과 명제 · 1. 집합 · 04 여집합과 차집합 · 2/2차시
\documentclass[11pt,a4paper]{article}
\usepackage{kotex}
\usepackage{iftex}
\ifPDFTeX\else
  \setmainhangulfont{Noto Serif CJK KR}
  \setsanshangulfont{Noto Sans CJK KR}
\fi
\usepackage[margin=2.1cm]{geometry}
\usepackage{parskip}
\usepackage{enumitem}
\usepackage{array}
\usepackage{tabularx}
\usepackage{booktabs}
\usepackage{xcolor}
\usepackage{amssymb}
\usepackage{tikz}
\usepackage[most]{tcolorbox}
\usepackage{fancyhdr}
\usepackage{titlesec}

\definecolor{goldc}{HTML}{B07C22}
\definecolor{goldstrong}{HTML}{8A5F16}
\definecolor{indigoc}{HTML}{2B3B57}
\definecolor{indigosoft}{HTML}{6E82A6}
\definecolor{linec}{HTML}{D6DACB}
\definecolor{surface2}{HTML}{E4E8DC}
\definecolor{notebg}{HTML}{FBF3E3}
\definecolor{inksoft}{HTML}{52604F}

\titleformat{\section}{\normalfont\sffamily\bfseries\large\color{indigoc}}{\thesection}{0.6em}{}
\titlespacing{\section}{0pt}{16pt}{8pt}
\newtcolorbox{ideabox}{enhanced, breakable, colback=white, colframe=goldc, boxrule=0.5pt, leftrule=3pt, arc=2pt, left=10pt, right=10pt, top=8pt, bottom=8pt}
\newtcolorbox{calloutbox}{enhanced, breakable, colback=white, colframe=indigosoft, boxrule=0.5pt, leftrule=3pt, arc=2pt, left=10pt, right=10pt, top=7pt, bottom=7pt}
\newtcolorbox{notebox}{enhanced, breakable, colback=notebg, colframe=goldc, boxrule=0.5pt, leftrule=3pt, arc=2pt, left=10pt, right=10pt, top=7pt, bottom=7pt}
\newtcolorbox{answerbox}{enhanced, breakable, colback=white, colframe=linec, boxrule=0.5pt, arc=2pt, left=10pt, right=10pt, top=7pt, bottom=7pt}
\newcommand{\lessonbanner}[3]{%
  \begin{tcolorbox}[enhanced, colback=indigoc, colframe=indigoc, boxrule=0pt, arc=0pt,
    left=14pt, right=14pt, top=14pt, bottom=10pt, borderline south={2.4pt}{0pt}{goldc}, fontupper=\color{white}]
    {\sffamily\footnotesize\color{white!85!black} #1}\\[4pt]{\sffamily\bfseries\Large #2}\\[3pt]{\sffamily\small\color{white!88!black} #3}
  \end{tcolorbox}}
\pagestyle{fancy}\fancyhf{}\renewcommand{\headrulewidth}{0pt}
\fancyfoot[L]{\footnotesize\color{inksoft} 공통수학2 · 2026학년도 2학기 · 지평선고등학교}
\fancyfoot[R]{\footnotesize\color{inksoft} 교사용 지도안 -- 배포 금지}

\begin{document}
\lessonbanner{공통수학2 · Ⅱ. 집합과 명제 · 1. 집합}{04 여집합과 차집합 --- 2/2차시}{드모르간의 법칙과 실생활 활용 · 교사용 지도안 (2026학년도 2학기)}
\vspace{10pt}

\noindent\begin{tabularx}{\linewidth}{@{}>{\bfseries\color{indigoc}}p{2.6cm}X@{}}
\toprule
대상 & 고1 · 2개 반 · 반당 13명 \\
차시 & 50분 (도입 10 · 전개 30 · 정리 10) \\
성취기준 & 집합의 연산을 수행하고, 벤 다이어그램을 이용하여 나타낼 수 있다. \\
선행 학습 & 23차시 --- 여집합과 차집합의 뜻 \\
\bottomrule
\end{tabularx}

\section{학습 목표 \& Big Idea}
\begin{ideabox}
\textbf{\color{goldstrong}영속적 이해 (Big Idea)}\\
`합집합의 여집합'은 `여집합들의 교집합'과 정확히 같다(드모르간의 법칙) --- `또는'을 뒤집으면 `그리고'가 된다는, 논리와 집합을 잇는 놀라운 대칭.
\vspace{4pt}
\small\color{inksoft} 핵심개념: \textbf{드모르간의 법칙} \quad\textbullet\quad 핵심개념: \textbf{논리와 집합의 대응} \quad\textbullet\quad 일반화: \textbf{`아니다'를 분배하면 `또는'과 `그리고'가 뒤바뀐다}
\end{ideabox}
\begin{calloutbox}
\textbf{차시 학습 목표}
\begin{itemize}[leftmargin=1.4em, itemsep=2pt, topsep=4pt]
  \item 여집합과 차집합의 성질을 이해하고 활용할 수 있다.
  \item 드모르간의 법칙을 이해하고, 이를 이용해 원소의 개수를 구할 수 있다.
\end{itemize}
\end{calloutbox}

\section{질문의 연결}
\begin{tabularx}{\linewidth}{@{}>{\bfseries\color{indigoc}\small}p{2.4cm}X@{}}
사실적 질문 & $(A\cup B)^C$와 $A^C\cap B^C$는 왜 같은 집합일까? \\[6pt]
개념적 질문 & `$A$이거나 $B$가 아니다'는 것이 `$A$도 아니고 $B$도 아니다'와 같다는 것을, 일상 언어의 예로 설명할 수 있을까? \\[6pt]
논쟁적 질문 & 두 신청(직업체험·대학탐방) 중 `아무것도 신청하지 않은 사람'의 수를 아는 것이 왜 전체를 계획하는 데 중요할까? \\
\end{tabularx}

\section{수업 흐름}
\begin{tabularx}{\linewidth}{@{}>{\centering\arraybackslash}p{1.6cm}X>{\raggedright\arraybackslash\small}p{3.1cm}@{}}
\toprule
\textbf{단계} & \textbf{활동 \& 발문} & \textbf{ATL / VTR} \\
\midrule
\textcolor{goldstrong}{\textbf{도입}}\newline\footnotesize 10분 &
\textbf{마음공부 (2분)} --- 유무념 대조.\newline
\textbf{생각 열기 (8분)} --- 여집합·차집합 성질 정리($U^C=\varnothing$, $\varnothing^C=U$, $(A^C)^C=A$, $A\cup A^C=U$, $A\cap A^C=\varnothing$, $A-B=A\cap B^C$)를 빠르게 복습. &
ATL: 자기관리(복습)\newline VTR: Connect-Extend-Challenge \\
\addlinespace
\textcolor{goldstrong}{\textbf{전개}}\newline\footnotesize 30분 &
\textbf{드모르간의 법칙 (15분)} --- 함께하기(벤 다이어그램 색칠로 $(A\cup B)^C=A^C\cap B^C$ 확인) $\to$ 법칙 정리 $\to$ 문제 5(교집합 버전).\newline
\textbf{적용 (15분)} --- 예제 2(원소 개수 계산) 함께 풀이 $\to$ 문제 7 짝 활동 $\to$ 문제 8(학급 설문 실생활 문제) 모둠 활동. &
ATT: 촉진자 역할\newline ATL: 협업·적용 \\
\addlinespace
\textcolor{goldstrong}{\textbf{정리}}\newline\footnotesize 10분 &
\textbf{개념 연결 질문 (5분)} --- ``드모르간의 법칙에서 $\cup$과 $\cap$이 어떻게 바뀌는지 한 문장으로 설명하면?''\newline
\textbf{성찰 (5분)} --- 감사 일기 1문장 + `집합' 대단원 마무리 + 다음 차시 예고(중단원 마무리). &
마음공부 · 감사 일기 \\
\bottomrule
\end{tabularx}

\begin{calloutbox}
\textbf{지도상의 유의점} --- 드모르간의 법칙 $(A\cup B)^C=A^C\cap B^C$, $(A\cap B)^C=A^C\cup B^C$에서 $\cup\leftrightarrow\cap$이 서로 바뀐다는 점을 명확히 강조한다. 문제 8과 같은 실생활 문제는 전체집합 $U$, 조건에 맞는 부분집합을 학생이 스스로 설정하는 연습이 핵심이다.
\end{calloutbox}

\section{형성평가 \& 답안}
\begin{answerbox}
\textbf{예제 2} \quad $n(U)=40$, $n(A)=21$, $n(B)=17$, $n(A\cup B)=34$일 때 $n(A^C\cap B^C)$\\
\textbf{\color{goldstrong}답} \; $n(A\cap B)=21+17-34=4$ $\to$ $n(A^C\cap B^C)=n((A\cup B)^C)=40-34=6$
\end{answerbox}
\begin{minipage}[t]{0.48\linewidth}
\begin{answerbox}
\textbf{문제 7} \quad $n(U)=50$, $n(A)=30$, $n(B)=21$, $n(A\cap B)=10$일 때 $n(A^C\cap B^C)$\\[4pt]
\textbf{\color{goldstrong}답} \; $n(A\cup B)=41$, $n(A^C\cap B^C)=50-41=9$
\end{answerbox}
\end{minipage}\hfill
\begin{minipage}[t]{0.48\linewidth}
\begin{answerbox}
\textbf{문제 8} \quad 학생 25명, 둘 다 신청 6명, 둘 다 안 함 2명, 직업체험이 대학탐방보다 5명 많음\\[4pt]
\textbf{\color{goldstrong}답} \; ⑴ $n(A\cup B)=23$ \quad ⑵ $n(A)=17$
\end{answerbox}
\end{minipage}

\section{성취수준 (이 차시 범위)}
\begin{tabularx}{\linewidth}{@{}>{\bfseries\color{goldstrong}}p{0.9cm}X@{}}
\toprule
A & 여집합·차집합의 연산을 체계적으로 수행하고, 연산 법칙을 벤 다이어그램으로 확인할 수 있다. \\
B & 여집합·차집합의 연산을 수행하고, 연산 법칙을 안다. \\
C & 여집합·차집합의 연산을 수행할 수 있다. \\
D & 간단한 여집합·차집합의 연산을 수행할 수 있다. \\
E & 간단한 두 집합의 여집합·차집합 연산을 수행할 수 있다. \\
\bottomrule
\end{tabularx}

\section{다음 차시}
\begin{calloutbox}
\textbf{중단원 마무리 --- 1. 집합 (25차시).} `여집합과 차집합'이 여기서 마무리되어 `1. 집합' 중단원 전체(18$\sim$24차시)가 끝난다. 다음 시간은 종합 정리 및 형성평가.
\end{calloutbox}
\end{document}
